% partinfo 章：Climate Change Adaptation 领域的研究脉络（活页章）
% 建档/更新日期由章首 \chdate 渲染（grep: chdate）
% 维护纪律：每落地或大改一篇精读章，必须回到本章补定位、改叙事
%           （CONTRIBUTING「partinfo 维护纪律」）。
% 体例注记: 2026-10-21 起，正文中的 terms/expressions 一律用英文原词
%          （weather、climate、location、panel approach……），
%          中文仅保留句法成分；定义处随文给中文对照并登记英文 \index。
% 方法线核对（bjr 2026-10-18 提出，2026-11-04 扩展）：方法线按正典
%          papers/01_methods/ 下的原文逐站核对——deschenes2007economic、
%          mcintosh2006identifying（工作稿）、burke2015global、
%          hsiang2016climate（所据 NBER w22181 工作稿版）、merel2021climate，
%          关键主张照引原句（CONTRIBUTING「原句引用」2026-10-29 条）；
%          后续改方法线仍须回各正典 PDF 核对页码与版本。
% 方法线重构（2026-11-19）：按主笔 2026-11-18 的修改意见，方法线由原来的五站
%          单线改为两条线——cross-sectional line（cross-sectional approach →
%          long differences → multistage model）与 panel line（panel approach →
%          nonlinear panel approach → climate penalty approach），Hsiang (2016)
%          作为把两条线放进同一框架的一站单列。原五站的原句引用全部保留并逐条
%          对照本地 PDF 改回原文字面；新增 burke2016adaptation、da2024climate
%          两篇的原句各若干。multistage model 一站所据的 Butler 与 Huybers（2013）
%          与 Kolstad 与 Moore 的综述（所据 NBER w25537 工作稿版）于 2026-11-19 下载
%          入 papers/01_methods/，该站已据原文写出。

\chapter{Climate Change Adaptation 研究脉络梳理}
\label{ch:info-climate-change-adaptation}
\chdate{2026-10-18}{2026-11-19}

本 part 梳理 climate change adaptation 研究的脉络。领域围绕的问题只有一个：climate change 会带来多大损害，经济主体又能通过长期 adaptation 抵消其中多少。本领域随 climate change 议题的兴起而出现；随着 adaptation 成为研究重点，方法重心从「估计影响」移向「检验 adaptation」。发展史按「当时面对什么问题、用什么方法回答、得到了什么发现、遗留了什么问题」登记，被推翻或修正的结论照实收录。


本 part 的两个术语区分在这里先立起来（identification 上的含义随精读章展开）。第一，\index{weather}weather（天气）指年际与年内的波动，\index{climate}climate（气候）指 weather 的长期分布：用 weather 波动识别的是 short-run response（不含 long-run adaptation），long-run adaptation 要靠横截面 climatic differences 或 long differences 这类长期 variation（变异）识别，凡引用识别结论先分清 variation 来源。第二，\index{adaptation}adaptation（适应）指长期发生的调整（品种、灌溉、种植制度），\index{coping}coping（应对）指短期风险管理（借贷、储蓄、转移），行文不混用。

\section{时间线}

以年份为轴登记方法演进与代表性发现；写「待补」的格子是欠账，对应精读章落地即还（自查：\texttt{grep -n '待补' part/*/partinfo.tex}）。

\begin{itemize}
  \item \textbf{1994}：Mendelsohn、Nordhaus 与 Shaw 用美国县横截面上的农地价值对 climate 变量回归，开截面 \index{Ricardian model}Ricardian 路线 \citep{mendelsohn1994impact}。
  \item \textbf{2006}：McIntosh 与 Schlenker 的工作稿给出含 penalty 结构 $(x_{it}-\mu_i)^2$ 的 DGP 与「hybrid estimator」（全局平方项与组内平方项并设），并指出二次项使 group means 重新进入识别：``a quadratic explanatory variable causes group means to re-enter the identification''\citep[摘要页；模型见第~4 页、hybrid estimator 见第~5--6 页]{mcintosh2006identifying}。\chref{ch:merel-gammans-2021} 的模型、增广回归与 proposition 皆源于此。
  \item \textbf{2007}：Deschênes 与 Greenstone 把 weather 波动面板（panel approach）引入农业影响评估（county fixed effects 加年际 weather 波动），回应截面路线的 omitted variables 批评 \citep{deschenes2007economic}（\chref{ch:deschenes-greenstone-2007}）。
  \item \textbf{2013}：Butler 与 Huybers 逐县估计美国玉米单产对高温的敏感度，再用敏感度在各县之间的差异刻画 adaptation，是 multistage model（多阶段模型）的首例 \citep{butler2013adaptation}。
  \item \textbf{2015}：Burke、Hsiang 与 Miguel 用 166 国 1960--2010 年度面板估计温度对经济产出的 nonlinear effect，响应在年均 13~$^\circ$C 处最优 \citep{burke2015global}（\chref{ch:burke-hsiang-miguel-2015}）；二次设定是面板文献的常见选择（名单见 \citealp[第~1208 页脚注~3]{merel2021climate}，其中既有这篇论文，也有 Deschênes 与 Greenstone 2007）。
  \item \textbf{2016}：Burke 与 Emerick 用 \index{long differences}long differences（两个长时段均值之差），对比 long-run response 与 panel approach 给出的 short-run response，考察美国农业的 adaptation \citep{burke2016adaptation}。同年 Hsiang 的综述系统整理了 climate econometrics 的各种 research designs \citep{hsiang2016climate}（\chref{ch:hsiang-2016}）。
  \item \textbf{2021}：Mérel 与 Gammans 把「nonlinear panel approach 的系数在多大程度上反映 long-run adaptation」形式化为可计算的 asymptotic bias 分解问题，并提出 climate penalty 增广回归（\chref{ch:merel-gammans-2021}）。
  \item \textbf{2024}：Da、Xu、Yi 与 McCarl 把 climate penalty model 用于中国冬小麦县级面板，配合分季节的 degree-days 变量 \citep{da2024climate}（\chref{ch:da-et-al-2024}）。
  \item 待补：Schlenker 与 Roberts（2009）的非线性温度响应一站；结构外推一站；各点落位后本时间线与「发现线」互相对照修订。
\end{itemize}

\section{方法线}

方法沿两条线演进，划分的依据是用来识别 climate 效应的 variation 来自哪里。\textbf{cross-sectional line} 用 location 之间的 climatic differences：从最早的 cross-sectional approach，到 long differences 与 multistage model。\textbf{panel line} 用 location 内部逐年的 weather fluctuation：从最早的 panel approach，经 nonlinear panel approach，到 climate penalty approach。两条线面对的是同一组取舍：横截面 variation 含有 long-run adaptation，但容易混入 omitted variables；组内 variation 干净，但测到的是否含 adaptation 存疑。每一站都按「上一站解决不了什么问题，这一站用什么 variation 回答」登记；各站的关键主张尽量照引原文原句（CONTRIBUTING「原句引用」），页码为所据版本页码（Hsiang 2016 所据为 NBER 工作稿版，见本章参考文献注）。

\subsection{Cross-sectional line：用 location 之间的 climatic differences}

\textbf{第一站：cross-sectional approach（Ricardian，1994，见时间线）。}用横截面上的农地价值对 climate 变量回归。\index{cross-sectional approach}Cross-sectional approach（截面法）的长处，\citet{hsiang2016climate} 的概括是：长期居住于一地的农民深知本地 climate 并据此调整投资与管理，``farm prices can be assumed to reflect all direct effects as well as all belief effects''\citep[第~5 页；\index{direct effect}direct effects（直接效应）指 climate 经由 weather 实现的直接影响，\index{belief effect}belief effects（信念效应）指主体因相信本地 climate 如此而做的调整，即 adaptation 行为，区分随 Hsiang 2016]{hsiang2016climate}。其弱点是 omitted variables：``A weakness of the cross-sectional approach is its vulnerability to omitted variables bias''\citep[第~5 页]{hsiang2016climate}，且 ``an econometrician can never be certain their model is unbiased''\citep[第~5 页]{hsiang2016climate}。\citet{deschenes2007economic} 用数据证实了这一弱点：同一 hedonic 设定下，估计值对控制变量、样本与加权的微小选择极其敏感，``it produces estimates that are extremely sensitive to seemingly minor choices about control variables, sample, and weighting''\citep[第~354 页]{deschenes2007economic}。这一站留下的问题是：能不能既用到 climate 的长期差异，又摆脱 location 之间不可观测的水平差异。

\textbf{第二站：long differences（2016，见时间线）。}\citet{burke2016adaptation} 的做法是对每个县取两个长时段的均值之差，再跨县比较这些长期变化：``We adopt a `long differences' approach and model county-level changes in agricultural outcomes over time as a function of these changes in temperature and precipitation, accounting for time-invariant unobservables at the county level and time-trending unobservables at the state level''\citep[第~107 页]{burke2016adaptation}。它保留了横截面比较（各县变暖的幅度不同），但比较的是变化而不是水平，县级不随时间变化的不可观测因素因此被差分掉。\citet{hsiang2016climate} 给它的方法定位是折中：``An approach that aims to compromise between the strengths and weaknesses of cross-sectional analysis and time series identification is the long differences strategy, where changes for both the outcome and the climate within locations are correlated across locations''\citep[第~7 页]{hsiang2016climate}。long differences 给出的 long-run 估计可以与 panel approach 的估计对比，两者之差被读作 adaptation；Burke 与 Emerick 的结论见发现线。代价有两条：识别所用的是各 location 长期趋势之间的差异，样本里必须有足够多的 location 和足够大的趋势差异；作对比基准的面板估计本身是否纯粹是 short-run response，要到 panel line 第三站才有答案。

\textbf{第三站：\index{multistage model}multistage model（多阶段模型；Butler 与 Huybers 2013）。}这一站把「响应本身随 climate 而变」直接估出来，分两步做。\citet{kolstad2019estimating} 的综述这样概括：``A first step models the local-linear effect of weather variation, allowing this linear response to weather fluctuations to vary across space, all linearly. In a second step, the coefficient on weather (from the first step) is used as the left hand side variable with a nonlinear function of climate and other control variables on the right hand side''\citep[第~11 页，NBER w25537 版页码]{kolstad2019estimating}。第一步用各 location 内部的 weather 波动，第二步用 location 之间的 climatic differences。Kolstad 与 Moore 因此把它与 long differences 一并归入兼用两种 variation 的 hybrid approaches（第~6 页）；本讲义随 \citet{da2024climate} 把它放在 cross-sectional line 上，因为 long-run adaptation 的信息全部来自第二步的横截面比较。首例是 \citet{butler2013adaptation}：美国东部 1{,}600 多个县、1981--2008 年的玉米单产，温度的影响用 growing degree days 与 killing degree days（日最高温超过 29\,$^\circ$C 部分的累积）刻画；各地单产对后者的敏感度不同，较热的地方敏感度较低。把这种空间上的 adaptation 当作未来 adaptation 的替代，``we find that losses to average US maize yields from a 2~$^\circ$C warming would be reduced from 14\% to only 6\% and that loss in net production is wholly averted''\citep[第~68 页]{butler2013adaptation}。这一站的代价有两条。其一，第二步是截面回归，第一站的 omitted variables 问题在这里重现：``any variable correlated with climate that also affects the marginal response to weather will bias the estimate of the effect of climate if not explicitly controlled for''\citep[第~12 页]{kolstad2019estimating}。其二，它与 long differences 一样需要足够多的 location，\citet{da2024climate} 的说法是：``These approaches exploit cross-sectional variation and thus require a large number of counties to obtain reliable estimates''\citep[第~825 页]{da2024climate}。Butler 与 Huybers 一文本讲义只读了与设计有关的部分，未设精读章。

\subsection{Panel line：用 location 内部的 weather fluctuation}

\textbf{第一站：\index{panel approach}panel approach（面板法）的引入（2007，\chref{ch:deschenes-greenstone-2007}）。}\citet{deschenes2007economic} 的设计陈述是 ``This paper's idea is simple—we exploit the presumably random year-to-year variation in temperature and precipitation to estimate whether agricultural profits are higher or lower in years that are warmer and wetter''\citep[第~355 页]{deschenes2007economic}，估计 ``These data are used to estimate the effect of weather on agricultural profits and yields, conditional on county and state by year fixed effects''\citep[第~355 页；此思路本身可追溯到 Huntington 1922，考据见 \citealp{hsiang2016climate} 第~6--7 页]{deschenes2007economic}。它回应的正是 cross-sectional line 第一站的 omitted variables 批评：``The core benefit of this approach is that it accounts for unobservable differences between units, eliminating a potential source of omitted variables bias''\citep[第~7 页]{hsiang2016climate}。variation 来源换成逐年 \index{weather fluctuation}weather fluctuation（天气波动）后，两重局限随即显形，且都是 D-G 自己写明的。其一，climate 信息进不了模型：``Second, it is impossible to estimate the effect of the long-run climate averages in a model with county fixed effects, because there is no temporal variation in $\overline{W}_{ic}$''\citep[第~367 页；$\overline{W}_{ic}$ 为 climate normals，即各县的长期 climate 变量——\citet{merel2021climate} 第~1207 页所引即此句]{deschenes2007economic}。其二，面板估计到的是不含全部 substitution 的响应：``The import for the subsequent analysis is that our estimates of the impact of climate change may be downward biased, relative to the preferred long-run effect that allows for all economic substitutions''\citep[第~359 页]{deschenes2007economic}——这句自认是后来 adaptation 批评的源头。

\textbf{第二站：\index{nonlinear panel approach}nonlinear panel approach（非线性面板法，weather 变量以二次项等非线性形式进入面板）。}这一站的模型源头是 \citet{mcintosh2006identifying}（2006，工作稿）：其式 (1) 即含 penalty 结构 $(x_{it}-\mu_i)^2$ 的 DGP，其 ``hybrid estimator''（全局平方项与组内平方项并设）是 \chref{ch:merel-gammans-2021} 增广回归的前身，摘要一句 ``a quadratic explanatory variable causes group means to re-enter the identification''\citep[摘要页；式 (1) 在第~4 页，hybrid estimator 在第~5--6 页]{mcintosh2006identifying} 即 proposition 的最早表述；``climate penalty'' 这个名字则由 Mérel 与 Gammans 命名\citep[第~1213 页]{merel2021climate}。宏观尺度上的代表应用是 \chref{ch:burke-hsiang-miguel-2015}：\citet{burke2015global} 用 166 国 1960--2010 年度面板，``comparing a country to itself in years when it is exposed to warmer- versus cooler-than-average temperatures due to naturally occurring stochastic atmospheric changes''\citep[第~236 页]{burke2015global}，配二次温度项，得到 ``country-level economic production is smooth, non-linear, and concave in temperature (Fig.~2a), with a maximum at 13~$^\circ$C''\citep[第~236 页]{burke2015global}。这一步改变了 identification 的 variation 构成——非线性项与 \index{fixed effects}fixed effects 并用，使横截面 \index{climatic differences}climatic differences（气候差异）重新进入模型识别——由此形成 \citet{merel2021climate} 整理的 proposition（命题）：``Including nonlinearities in explanatory weather variables in a fixed-effects panel regression allows cross-sectional weather variation to enter model identification, suggesting that the estimates obtained from a nonlinear panel model at least partially reflect long-run adaptation to climate''\citep[第~1208 页]{merel2021climate}。该 proposition 在 2021 年前已被多项研究谨慎承认（名单见同页正文，其中有 \citealp{mcintosh2006identifying} 与 \citealp{burke2015global}），但它只回答了「估计里有没有 long-run 成分」：``the extent to which nonlinear panel estimates should be thought of as inclusive of climatic adaptation remains unclear''\citep[第~1208 页]{merel2021climate}。两点须防误读：其一，\citet{burke2015global} 是这一站的旗舰应用而非首创；其二，「线性对非线性」是教学简化——\citet{deschenes2007economic} 的利润方程已对 degree-days 与降水取二次（第~367 页），故它同列于 \citet{merel2021climate} 的二次设定面板名单（第~1208 页脚注~3）；真正决定系数如何解释的，是 penalty 结构是否在回归里。

\textbf{第三站：climate penalty approach（2021，\chref{ch:merel-gammans-2021}）。}这一站回答 proposition 悬置的「多少」：在 McIntosh 与 Schlenker（2006）提出的模型框架上（\citealp{mcintosh2006identifying}），把「nonlinear panel estimates 里 long-run 成分占多少」形式化为可计算的 \index{asymptotic bias}asymptotic bias（渐近偏差）分解。其一，``panel estimates from the quadratic weather-outcome relationship can be thought of as a weighted average of short- and long-run responses, with the weight on the long-run coefficients increasing with the share of overall weather variation attributable to cross-sectional (i.e., climatic) differences''\citep[第~1209 页]{merel2021climate}；weather 分布无 skewness 时为 convex combination（凸组合，第~1217 页），线性情形下偏差方向由条件于 location 的 weather 分布的 \index{skewness}skewness（偏度）决定，可正可负（第~1215 页）。其二，同一框架反向给出一致估计的路线：在面板中显式加入 \index{climate penalty}climate penalty term，一次回归同时得到 conditional on climate 的 short-run response、全局 long-run response 与一个 \index{climatic adaptation}climatic adaptation 的正式检验（第~1210、1220--1221 页）；该检验与 long differences 的检验同型，``our test of climatic adaptation is implicit in that it relies on the comparison of short-run versus long-run outcome responses''\citep[第~1210 页]{merel2021climate}。这一设定此后被直接搬用到新的作物与国家：本 part 的 Application 1（\chref{ch:da-et-al-2024}）把它用于中国冬小麦，配合分季节的 degree-days 变量，并报告了 penalty term 全部加入时出现的反常符号；它成为这条线上通行的做法。模型 terms 与解释 terms 在这条线上的逐条对应与显式分解，见精读章 \chref{ch:merel-gammans-2021}。

\subsection{把两条线放进同一框架：Hsiang（2016）}

\chref{ch:hsiang-2016} 把上述争论正式化。\citet{hsiang2016climate} 的出发点是 ``a debate as to whether regressions on `weather' variables provide meaningful insight into the effects of climate''\citep[第~2 页]{hsiang2016climate}，他给出 ``conditions under which the use of weather variables in regressions is justified and, perhaps surprisingly, dominates traditionally preferred methods''\citep[第~2 页]{hsiang2016climate}，并把领域整理成三种 research designs：``cross-sectional approaches, use of time-series variation, and a hybrid known as `long differences'''\citep[第~5 页]{hsiang2016climate}——前者与后者属 cross-sectional line，中间一种属 panel line。他提出两件工具。其一，\index{frequency-identification tradeoff}frequency-identification tradeoff（频率—识别权衡）：低频 variation 可能捕捉 adaptation，但识别可信度随之下降（第~8--9 页）；若 belief effects 主要通过改变 direct effects 起作用，``total effects of climate still may be nearly identified with high-frequency time-series''\citep[第~9 页]{hsiang2016climate}。其二，对 \index{marginal treatment comparability}marginal treatment comparability（边际处理可比性，第~6 页）的部分检验（第~9--10 页）：若三种设计的估计相等（$\hat\beta_{CS}=\hat\beta_{LD}=\hat\beta_{TS}$），weather 的边际变化即可当作 climate 边际变化的替代；他在玉米单产数据上实施的结果，除横截面外各估计统计上无异（第~10 页）。他还证明了一个正式结果：若结果变量已被优化，在边际处 weather 效应等于 climate 效应（\S 2.4）。与这一结果对照，Mérel 与 Gammans 自陈：``Whereas previous research has already argued that short-run and long-run responses should be identical at the margin whenever the outcome variable is optimized (e.g., Hsiang 2016), our result is both different and more specific''\citep[第~1217 页]{merel2021climate}——面板估计到的是加权平均，恰在边际处等同只是特例。

两条线各站的适用边界比较，以及与之平行的结构外推一站，待补。

\section{发现线}

待补：本领域已确认的重点结论与量级区间，随各精读章「结果与解释」节落地回填；数字必须带表号（页码）与样本口径。不同研究的 climate 变量定义、样本与被解释变量不同，下列数字不得跨条比较大小。目前已可登记六条。

其一与其二出自 \chref{ch:merel-gammans-2021}：摘要报告对美国与法国的作物单产找到了 long-run climatic adaptation 的一些证据；方法线悬置的「面板系数里 long-run 成分占多少」第一次有了量级回答——growing season 平均温度的效应应主要按 long-run response 解读，三个面板的长期权重分别为 98\%、88\%、86\%，growing season 降水的效应则含相当大的 short-run 成分，长期权重分别为 41\%、7\%、53\%（Table~2 与正文，第~1225 页；样本口径：美国县面板为 100 度经线以东、不含 Florida 的县，$I=2{,}391$；Corn Belt 三州 Iowa、Illinois、Indiana，$I=293$；法国 department 级面板，$I=88$；三面板均 $T=66$；美国的产出为玉米与大豆单产，法国为小麦与大麦单产，第~1224 页）。

其三出自 \chref{ch:deschenes-greenstone-2007}：``Using long-run climate change predictions, the preferred estimates indicate that climate change will lead to a \$1.3 billion (2002\$) or 4.0-percent increase in annual agricultural sector profits''\citep[第~356 页，AER 97(1) 已刊版口径]{deschenes2007economic}，且响应面在预测变化区间内平坦——``all of the response surfaces are flat over the ranges of the predicted changes in degree-days and inches''\citep[第~355 页]{deschenes2007economic}。

其四出自 \chref{ch:burke-hsiang-miguel-2015}：国家层面的经济产出对年平均温度呈光滑的凹响应，最优在 13~$^\circ$C（166 国、1960--2010 年、人均 GDP 增长率，第~236 页）；1960--1989 与 1990--2010 两段子样本的响应几乎相同，作者据此判断样本期内未见明显 adaptation（第~237 页）；基准 projection 下 2100 年全球产出比无 climate change 的情形低 23\%，影响为正的概率 0.29（SSP5，第~238 页）。

其五出自 Burke 与 Emerick（2016，本讲义未设精读章）：``Longer run adaptations appear to have mitigated less than half—and more likely none—of the large negative short-run impacts of extreme heat on productivity''\citep[第~106 页，摘要]{burke2016adaptation}（美国县级农业，long differences 对面板估计）。

其六出自 \chref{ch:da-et-al-2024}（中国黄淮海 352 县、1981--2015 年、冬小麦对数单产）：是否纳入 freezing 变量决定 projection 的符号，不计时各情景减产 12.5\% 到 48.2\%，计入时小幅增产 1.9\% 到 2.5\%（第~817、837 页）；long-run 口径下 Fall 高于 24\,$^\circ$C 的 degree-days 每多 10 个减产 7.0\%，conventional 面板的对应数字是 11.4\%（Table~2，第~829、832 页）。

待补：被推翻或修正的结论（如早期 Ricardian 外推的争议）。已可登记两例。第一例，\chref{ch:deschenes-greenstone-2007} 重做 hedonic 估计，``estimates of the effect of the benchmark doubling of greenhouse gases on the value of agricultural land range from $-\$200$ billion (2002\$) to \$320 billion (or $-18$ percent to 29 percent)''\citep[第~357 页]{deschenes2007economic}，比此前文献记录的区间更宽，结论是 ``Despite its theoretical appeal, we conclude that the hedonic method may be unreliable in this setting''\citep[第~357 页]{deschenes2007economic}。第二例是对一种做法的修正：把 long differences 估计与 conventional 面板估计之差当作 adaptation 的量度，前提是面板估计代表 short-run response；\chref{ch:merel-gammans-2021} 说明面板估计本身是 short-run 与 long-run response 的加权平均，\chref{ch:da-et-al-2024} 据此指出这种比较可能低估 adaptation（该文第~836 页）。

\section{章节速查表}

本表是「章与 partinfo 互指」的锚：每篇精读章一格。讲方法的各章在前，应用篇在后；应用篇按收入本 part 的先后编号（Application 1、2……），编号出现在目录、页眉、本表与各章的坐标句里，章题仍用论文标题。

\begin{table}[htbp]
\centering
\small
\begin{tabular}{@{}p{1.4cm}p{2.8cm}p{3.0cm}p{2.7cm}p{3.2cm}@{}}
\toprule
章 & 论文 & 回答的问题 & 方法 & 给领域留下了什么 \\
\midrule
\chref{ch:deschenes-greenstone-2007} & The Economic Impacts of Climate Change: Evidence from Agricultural Output and Random Fluctuations in Weather（AER 2007） & 年际 weather 波动识别的 climate change 对美国农业利润的影响有多大？hedonic 截面估计是否可靠？ & county 与 state$\times$year fixed effects 下的年际 weather 波动回归（panel line 第一站） & hedonic 估计对口径选择极敏感、不可靠的重估证据；自认面板估计不含全部 substitution、相对 long-run effect 向下偏（adaptation 批评的源头） \\
\chref{ch:burke-hsiang-miguel-2015} & Global Non-Linear Effect of Temperature on Economic Production（Nature 2015） & 温度与经济产出在全球国家面板上是否非线性？样本期内的变暖是否带来了 adaptation？ & 166 国 1960--2010 增长率面板，二次温度项（country 与 year fixed effects、国家二次趋势；panel line 第二站的旗舰应用） & 全球温度—产出响应在 13~$^\circ$C 处最优的非线性证据；「样本期内未见 adaptation」的判断引发后续 adaptation 检验 \\
\chref{ch:hsiang-2016} & Climate Econometrics（ARRE 2016） & weather 回归何时能识别 climate 效应？三种 research designs 如何取舍？ & 方法综述：正式化 weather／climate，把 cross-section、time-series、long differences 三种设计放进同一框架 & marginal treatment comparability、frequency-identification tradeoff；$\hat\beta_{CS}=\hat\beta_{LD}=\hat\beta_{TS}$ 部分检验（玉米单产应用）；边际处 weather 效应等于 climate 效应的正式结果 \\
\chref{ch:merel-gammans-2021} & Climate Econometrics: Can the Panel Approach Account for Long-Run Adaptation?（AJAE 2021） & nonlinear fixed effects 面板（忽略 penalty term 的二次设定）估计到的系数，在多大程度上能当作 long-run climate response 来读？ & 在 McIntosh 与 Schlenker（2006）框架上做 asymptotic bias 分解（固定 location 数、大 $T$）；climate penalty 增广回归同时估计 long-run 与 short-run response（panel line 第三站） & 线性情形偏差由条件于 location 的 weather skewness 驱动、可正可负；无 skewness 时面板系数是 short-run 与 long-run response 的 convex combination；空间大面板的温度系数已主要反映 long-run response（降水不然） \\
\chref{ch:da-et-al-2024}（Application 1） & Climate Change and Winter Wheat Yields: The Roles of Freezing Days and Long-Run Adaptations（ERAE 2024） & 中国冬小麦的 climate 影响评估是否必须分季节并纳入 freezing？climate penalty model 能否用于区分 short-run 与 long-run 影响？ & 黄淮海 352 县 1981--2015 面板；按季节由数据选阈值构造 degree-days；对两个变量加 penalty term（panel line 第三站的应用） & freezing 变量决定 projection 的符号；penalty term 全部加入时出现反常符号的实例；long differences 与面板之差可能低估 adaptation 的论点 \\
\bottomrule
\end{tabular}
\end{table}

% 章末参考文献（只列本章引用的条目），勿删
\chapterrefs
